% Lösungen Einheit 5 von 5 – Drei Variablen und Lösungsvielfalt (Sek II) (zusammenbau v0.1)
\einheitenkopf{Einheit 5 von 5 · Drei Variablen und Lösungsvielfalt (Sek II)}

\erg{29}{a) unendlich viele \quad b) I: $2 + 2 \cdot 1 - 1 = 3$ (wA); II: $3 \cdot 2 + 1 - 2 \cdot 1 = 5$ (wA). Also ist $(2;\,1;\,1)$ eine Lösung. \quad c) $z = 3$, $y = 2$, $x = 3$ \quad d) II $-$ III: $2z = 4$, $z = 2$; $y = 3$; $x = 3$ \quad e) I $+$ II ergibt $2x + y + 3z = 5$, das ist III: III bringt nichts Neues, das System hat unendlich viele Lösungen. \quad f) III ist das Dreifache von II. $z = t$, $x = 1 + t$, $y = 5 - t$; alle positiv: $0 < t < 5$; größtes $x$ bei $t = 4$: $x = 5$, $y = 1$, $z = 4$. \quad g) I und II: $x = 3$, $y = 2$; in III: $a = 2 \cdot 3 + 2 = 8$. Für $a = 8$ genau eine Lösung (nicht unendlich viele), für jedes andere $a$ keine. \quad h) $a^2 - 9 = (a - 3)(a + 3)$. $a = 3$: $0 = 6$ (fA), keine Lösung; $a = -3$: $0 = 0$ (wA), unendlich viele; sonst genau eine (z aus III, dann y aus II, x aus I). \quad i) $z \ge 2x$ heißt $t \ge 6 - 8t$, also $t \ge \frac{2}{3}$ (Grenze $t = \frac{2}{3}$). Weil $y$ ein Anteil ist, gilt $y = -2 + 3t \ge 0$ heißt ebenfalls $t \ge \frac{2}{3}$ – jede zulässige Lösung erfüllt die Aussage.}
\erg{30}{a) Für $a = 2$ wird III zu $0 = 6$ (fA): dann gibt es keine Lösung. Nur für $a \ne 2$ darf man durch $a - 2$ teilen. \quad b) $2 \cdot (-1) = -2$ wird mit $+2$ ausgeglichen: $y = 5 + 2 = 7$; dann $x = 2 - 7 + 1 = -4$. Lösung: $x = -4$, $y = 7$, $z = -1$. \quad c) $x = \frac{4}{3}$ ist nicht ganzzahlig; $t$ muss ein Vielfaches von $3$ sein. Richtig: $t = 3$, $x = 1$, $y = 2$, $z = 3$. \quad d) $0 \cdot z$ ist für jedes $z$ null. Ist $c = 0$, stimmt die Gleichung für jedes $z$ (unendlich viele Lösungen); ist $c \ne 0$, für keines (keine Lösung). \quad e) Zu jedem Wert von $t$ ergeben die übrigen Gleichungen passende $x$ und $y$ – jede Zahl $t$ gibt eine Lösung. Zusatzbedingungen werden zu Ungleichungen oder Teilbarkeitsbedingungen für $t$ und lassen nur bestimmte Werte übrig.}
